\chapter{Validação Técnica e Evidência em Tempo de Execução}

\section{Âmbito e Estratégia de Validação}

Este capítulo apresenta a validação técnica do NatureProtector V1. O objetivo é demonstrar que a baseline implementada executa a pipeline de prevenção de ponta a ponta, persiste estado operacional, expõe projeções através da API e produz evidência auditável em base de dados e testes automatizados.

A validação é técnica e operacional. Não constitui validação científica do modelo de risco, não prova calibração dos scores contra ocorrências reais e não estabelece equivalência com sistemas oficiais ou índices de perigo de incêndio. A sua função é verificar se o sistema funciona como pipeline de software rastreável. Esta distinção é necessária porque índices e sistemas reconhecidos de perigo de incêndio, como o FWI, produtos do IPMA e plataformas europeias de monitorização, dependem de metodologias, dados e processos de validação próprios~\citep{nrcan_fwi_system,ipma_fwi,ipma_pir_rcm,ec_jrc_effis}.

A estratégia de validação combinou quatro fontes principais:

\begin{itemize}
    \item inspeção da infraestrutura e do plano de controlo;
    \item queries SQL sobre os schemas \texttt{control}, \texttt{pipeline} e \texttt{projection};
    \item comparação entre respostas da Backoffice API e projeções persistidas;
    \item execução de testes automatizados, incluindo testes específicos da política de alertas.
\end{itemize}

A principal evidência foi recolhida na fase C7, dedicada à execução end-to-end, recolha de evidência, sincronização documental e auditoria técnica final. Os artefactos relevantes incluem:

\begin{itemize}
    \item \texttt{v1-runtime-evidence-2026-05-14-0107.md};
    \item \texttt{v1-dotnet-test-2026-05-14-final.txt};
    \item \texttt{v1-test-list-2026-05-14.txt};
    \item \texttt{v1-alert-policy-tests-2026-05-14-final.txt};
    \item \texttt{c7-runtime-evidence-summary-2026-05-14.md}, quando presente como síntese consolidada.
\end{itemize}

Estes nomes devem corresponder aos ficheiros finais existentes no repositório. Caso algum artefacto tenha nome diferente, o texto deve ser ajustado antes da entrega.

\section{Evidência da Pipeline e das Projeções}

A validação confirmou que a infraestrutura local estava disponível, que o plano de controlo continha a configuração esperada e que a área piloto \texttt{proenca-a-nova}, sensores, células e cenários estavam presentes para execução.

A evidência da pipeline centrou-se no schema \texttt{pipeline}. A tabela \texttt{pipeline.event\_inbox} confirmou que eventos produzidos pelo simulador chegaram à fronteira durável da pipeline. A tabela \texttt{pipeline.processing\_attempts} confirmou que os eventos não foram apenas recebidos, mas processados através de tentativas registadas. As tabelas \texttt{pipeline.rejected\_events} e \texttt{pipeline.quarantined\_events} permitiram distinguir rejeições por validação de falhas históricas ou retry esgotado. Esta separação entre receção, tentativa, confirmação e efeito persistido é coerente com práticas de processamento assíncrono e idempotente em sistemas orientados a mensagens~\citep{rabbitmq_consumer_acknowledgements,richardson_idempotent_consumer}.

A evidência das projeções centrou-se no schema \texttt{projection}. As tabelas \texttt{projection.accepted\_reading\_log} e \texttt{projection.risk\_assessment\_log} confirmaram a existência de leituras aceites e avaliações de risco persistidas. As tabelas \texttt{projection.area\_risk\_snapshot\_log}, \texttt{projection.cell\_operational\_state} e \texttt{projection.area\_operational\_state} confirmaram que os resultados da pipeline foram agregados em modelos de leitura operacional. A tabela \texttt{projection.alert\_state} confirmou suporte persistido para estado de alerta.

\begin{table}[h]
\centering
\caption{Evidência técnica principal da C7.}
\label{tab:c7-main-evidence}
\begin{tabular}{p{0.32\textwidth} p{0.50\textwidth}}
\hline
\textbf{Evidência} & \textbf{Interpretação} \\
\hline
\texttt{event\_inbox} e \texttt{processing\_attempts} & Eventos chegaram à pipeline e foram processados através de tentativas registadas. \\
\texttt{rejected\_events} & Rejeições foram classificadas antes da pipeline de risco aceite. \\
\texttt{quarantined\_events} & Registos de quarentena foram identificados e datados. \\
\texttt{accepted\_reading\_log} e \texttt{risk\_assessment\_log} & Leituras aceites e avaliações de risco foram produzidas e persistidas. \\
\texttt{area\_operational\_state} e \texttt{cell\_operational\_state} & Projeções operacionais foram atualizadas por área e por célula. \\
Backoffice API & A API expôs estado coerente com projeções persistidas. \\
Testes automatizados & A suite passou e cobriu comportamentos críticos, incluindo alertas. \\
\hline
\end{tabular}
\end{table}

Esta evidência demonstra que a V1 não ficou apenas pela receção de eventos. Os eventos progrediram pela pipeline, produziram avaliações candidatas de risco e atualizaram projeções operacionais. Isto suporta validação técnica, mas não validação científica dos valores calculados.

\section{Rejeições, Quarentenas e Erros Históricos}

Eventos rejeitados e eventos em quarentena foram analisados separadamente, porque representam situações diferentes. Os eventos rejeitados foram classificados como \texttt{invalid\_operational\_state}. Isto significa que a pipeline recusou leituras com estado operacional inválido antes de estas entrarem na pipeline de risco aceite. Esta rejeição não indica, por si só, erro da implementação; demonstra que a fronteira de validação está ativa.

Os eventos em quarentena foram classificados como \texttt{retries\_exhausted}. A análise temporal mostrou que estes registos estavam associados a execuções anteriores, não à execução final da C7. Por isso, foram documentados como registos históricos de retry esgotado.

Durante a análise surgiu ainda o erro:

\begin{quote}
\texttt{The given key 'EmptyProjectionMember' was not present in the dictionary.}
\end{quote}

Este erro estava associado a registos \texttt{processing\_failed}. A análise temporal mostrou que os registos eram históricos e não afetavam a execução final da C7. O erro foi, portanto, classificado como histórico e não bloqueante.

Esta classificação é importante porque evita uma conclusão incorreta. Nem toda a rejeição indica falha, e nem todo o erro antigo representa problema ativo. A validação deve classificar, datar e interpretar os registos, não apenas contá-los.

\section{Consistência entre API, Alertas e Testes}

A Backoffice API foi validada por comparação com a projeção persistida em \texttt{projection.area\_operational\_state}. O endpoint de estado operacional expôs campos como código da área, versão de configuração, timestamp do snapshot, score agregado, nível de risco, severidade, resumo, número de avaliações, atualização e estado de alerta.

Na evidência final C7, o estado da área apresentava score agregado próximo de \texttt{0.448}, nível agregado \texttt{Moderate}, severidade \texttt{Medium} e \texttt{75} avaliações. A API expôs estado coerente com a projeção persistida. O endpoint de alertas ativos devolveu uma lista vazia, coerente com o score agregado estar abaixo do limiar de aviso. A ausência de alerta ativo neste cenário não constitui falha da política de alertas; apenas indica que o estado observado não cumpria condições de abertura.

A política de alertas foi validada sobretudo por testes automatizados, uma vez que o cenário final não atravessou necessariamente os limiares configurados. A lista de testes incluiu transições de \texttt{None} para \texttt{Warning} e \texttt{Alarm}, histerese de \texttt{Warning} e \texttt{Alarm}, resolução abaixo do limiar de fecho, criação de alertas em projeções PostgreSQL e descida de \texttt{Alarm} para \texttt{Warning} sem resolução indevida.

A execução específica dos testes de alert policy passou. A execução completa de \texttt{dotnet test} também terminou sem falhas, cobrindo projetos de domínio, prevenção, infraestrutura, hosts, integração e Backoffice API. O output reportou um aviso de dependência relacionado com \texttt{OpenTelemetry.Exporter.OpenTelemetryProtocol}. Este aviso não invalida a validação técnica, mas deve ser tratado como item de manutenção.

A formulação conservadora é:

\begin{quote}
A evidência C7 mostra consistência entre a projeção operacional persistida e a resposta da API. Não foi observada evidência de recálculo divergente na API durante a recolha em tempo de execução. A política de alertas foi suportada por testes automatizados, embora o cenário final não tenha exigido alerta ativo.
\end{quote}

\section{Limitações da Evidência}

A evidência C7 suporta um veredito técnico positivo, mas deve ser interpretada dentro dos seguintes limites:

\begin{itemize}
    \item os dados utilizados são simulados;
    \item pesos, limiares, classes e valores de histerese são parâmetros candidatos;
    \item a validação é técnica e em tempo de execução, não científica;
    \item o cenário final não apresentou alerta ativo;
    \item alguns conceitos internos podem não estar expostos como colunas persistidas dedicadas;
    \item erros históricos existem na base de dados, embora tenham sido classificados como não bloqueantes;
    \item o aviso de dependência deve ser tratado como item de manutenção;
    \item uma árvore de trabalho Git não limpa, se aplicável, limita a rastreabilidade da evidência final.
\end{itemize}

Estas limitações não anulam a validação técnica. Definem o alcance correto da evidência e impedem que os resultados sejam interpretados como calibração científica ou prontidão para produção. Esta formulação é consistente com a distinção feita ao longo do relatório entre validação técnica da pipeline e validação científica de modelos ou índices de perigo de incêndio~\citep{nrcan_fwi_system,ipma_fwi,ec_jrc_effis}.

\section{Veredito Técnico}

Com base na evidência recolhida e nos testes automatizados, o veredito adequado para a baseline V1 é:

\begin{quote}
\textbf{OK com limitações.}
\end{quote}

A evidência suporta a conclusão de que o NatureProtector V1 executa uma pipeline local end-to-end do módulo de prevenção. O sistema recebe eventos simulados, persiste-os de forma durável, regista tentativas de processamento, classifica rejeições e quarentenas, produz leituras aceites e avaliações de risco, atualiza projeções operacionais, expõe estado pela API e valida a política de alertas através de testes.

Nenhum bloqueador P0 foi confirmado na análise final. O erro \texttt{EmptyProjectionMember} foi classificado como histórico. Eventos em quarentena foram classificados como registos históricos de retry esgotado. Eventos rejeitados foram classificados como rejeições deliberadas de estado operacional inválido. A API mostrou consistência com a projeção persistida. A política de alertas foi coberta por testes automatizados. A suite completa de testes passou.

\begin{table}[h]
\centering
\caption{Veredito técnico final da C7.}
\label{tab:technical-verdict}
\begin{tabular}{p{0.34\textwidth} p{0.50\textwidth}}
\hline
\textbf{Item} & \textbf{Classificação} \\
\hline
Infraestrutura e plano de controlo & Operacionais na evidência local. \\
Pipeline & Operacional, com evidência durável. \\
Eventos rejeitados & Estado operacional inválido, rejeição deliberada. \\
Eventos em quarentena & Registos históricos de retry esgotado. \\
Erro \texttt{EmptyProjectionMember} & Histórico e não bloqueante. \\
Avaliações de risco e projeções & Produzidas e persistidas. \\
API & Coerente com projeções persistidas. \\
Política de alertas & Coberta por testes automatizados. \\
Suite de testes & Passou. \\
Validação científica & Não reivindicada. \\
Veredito final & OK com limitações. \\
\hline
\end{tabular}
\end{table}
Assim, a conclusão correta é que o NatureProtector V1 se encontra tecnicamente operacional e pode ser classificado como \textit{OK com limitações}.